\chapter{Risk Management: Kelly, Drawdown, and Metrics}
\label{ch:kelly_risk}
\chaptermark{Risk management: Kelly, drawdown, metrics}

You have a positive-edge strategy. How much of your capital should
you bet on each trade? Too little and you leave money on the table;
too much and a losing streak wipes you out. The \textbf{Kelly
criterion} \citep{kelly1956new} gives the unique fraction of
capital that maximises the \emph{long-run geometric growth rate}
of wealth. This chapter derives Kelly in the binary case,
generalises to log-normal returns, and then steps back for a
practitioner reality check: \textbf{Kelly is theoretically optimal
and almost no professional trades full Kelly}, for reasons that
matter more than the derivation itself.

The second half switches from \emph{sizing} to \emph{scoring}: how
do you tell a good run from a lucky one, and compare two
strategies? That is the province of performance metrics --- Sharpe,
Sortino, Calmar, information ratio --- and of drawdown arithmetic,
which contains the single most important identity in this chapter:
the asymmetric recovery formula. A 20\% loss needs 25\% to break
even; 50\% needs 100\%; 80\% needs 400\%. The whole reason one
prefers a Sharpe of 1.5 with a 10\% max drawdown to a Sharpe of
2.5 with a 60\% max drawdown is baked into that asymmetry.

Notation: $f$ is the Kelly fraction of wealth staked, $p$ is win
probability, $q=1-p$, $b$ is win size, $a$ is loss size, $\mu$ is
expected log-return, $\volat$ log-return standard deviation, $R$ a
per-period return, $\rrf$ the risk-free rate. Execution scheduling
(\cref{ch:almgren_chriss}) told you \emph{when} to trade; this
chapter says \emph{how much} to hold and how to score what you did.

\begin{backgroundbox}[Prerequisites]
From earlier in the book:
\begin{itemize}
  \item \Cref{ch:mean_reversion_ou}: log-returns, i.i.d.\ bet
  sequences, and the language of expected value, variance, and
  long-run sample averages. The $O$--$U$ half-life of
  \cref{ch:mean_reversion_ou} is \emph{not} used here, but the
  habit of modelling increments as mean plus noise is.
  \item \Cref{ch:momentum}: the volatility-scaled position-sizing
  rule of time-series momentum --- ``hold $\volat_{\text{target}}
  / \volat_t$ units of the instrument'' --- is a \emph{heuristic
  fractional-Kelly rule in disguise}. This chapter explains why.
  \item \Cref{ch:vol_surface}: $\volat$ denotes annualised
  (log-)volatility. Throughout this chapter $\mu$ is expected
  log-return per unit time, $\volat$ is the standard deviation of
  log-returns per unit time, and the risk-free rate is $\rrf$.
  \item \Cref{ch:almgren_chriss}: mean--variance thinking ---
  trade off an expected quantity against a variance. Kelly will
  turn out to be \emph{mean--variance optimality in a particular
  disguise} (the $\log$ disguise) rather than a competitor to it.
\end{itemize}
From outside the book:
\begin{itemize}
  \item Basic \textbf{expected utility theory}: if you have a
  utility function $U(w)$ over wealth, a ``rational'' bettor
  maximises $\E[U(W_{\text{final}})]$ rather than
  $\E[W_{\text{final}}]$ itself. Kelly corresponds to the choice
  $U(w) = \ln w$.
  \item \textbf{Logarithms and exponentials}: $\ln(1+x) \approx x
  - x^2/2$ for small $x$, $\frac{\mathrm d}{\mathrm df}\ln(1+fx) =
  x/(1+fx)$, and the fact that $\ln$ of a product is the sum of
  $\ln$'s.
  \item \textbf{First-order conditions}: set the derivative of a
  concave function to zero to find the maximum.
\end{itemize}
\end{backgroundbox}

\section{The Kelly criterion: binary case}
\label{sec:kelly_binary}

\subsection{Setup}

A biased coin lands heads with probability $p > 1/2$ and tails
with probability $q = 1 - p$. Stake a dollar on heads: you win
$\$b$ on heads, lose $\$a$ on tails, with $a, b > 0$. The unit
even-money bet is $a = b = 1$; a sports bettor quoted at $+150$
decimal odds plays $b = 1.5$, $a = 1$; a slot machine with $99\%$
payout has $pb < qa$ and should not be played.

You bet a fraction $f \in [0, 1]$ of current wealth each trial.
With $W_0 = 1$ and trial-$i$ per-unit payoff $R_i \in \{+b, -a\}$,
\begin{equation}
  W_N = W_0 \prod_{i=1}^{N} (1 + f R_i).
\end{equation}
Wealth \emph{multiplies} trial-by-trial: a gain on trial 5 grows
the base on which trials 6, 7, \ldots compound. This is the single
most important distinction in the chapter and the reason plain
$\E[W_N]$ is the \emph{wrong} objective.

\subsection{Why maximising expected wealth fails}
\label{sec:kelly_expected_wealth_trap}

What fraction $f$ maximises $\E[W_N]$? After one trial,
\begin{equation}
  \E[W_1] = p(1 + fb) + q(1 - fa) = 1 + f(pb - qa).
\end{equation}
With positive edge $pb - qa > 0$ this is strictly increasing in
$f$, with maximum at $f = 1$: bet everything. By compounding, the
same holds for $\E[W_N]$.

But $f = 1$ is catastrophic: you lose your entire bankroll the
first tails, which happens with probability $1 - p^N \to 1$.
Expected wealth grows at the maximum rate only because one lucky
path of probability $p^N$ contributes an exponentially large
amount and drags the mean upwards \citeThorp. With probability one
you are broke. $\E[W_N]$ cannot distinguish ``probably rich'' from
``almost certainly broke but very rarely stupendously rich''.

\begin{pitfallbox}[Do not maximise expected wealth]
A strategy that maximises $\E[W_N]$ on a multiplicative bet is
generically a strategy that guarantees ruin with probability one.
The expectation is dominated by the one path where every bet wins,
which is exponentially rare but exponentially large. Kelly's
insight is that you should instead maximise the \emph{typical}
path, which is captured by $\E[\ln W_N]$. For compounding bets,
``expected log'' is the right objective; ``expected level'' is a
trap.
\end{pitfallbox}

\subsection{The Kelly objective: expected log-wealth}

Switch the objective from $\E[W_N]$ to $\E[\ln W_N]$ --- equivalently,
adopt the log-utility $U(w) = \ln w$. Since $\ln W_N = \ln W_0 +
\sum_i \ln(1 + fR_i)$ turns the multiplicative wealth process into
a sum, its expectation is $N$ times a \emph{per-trial} growth rate:
\begin{equation}
  G(f) \;\defeq\; \E[\ln(1 + fR)]
  \;=\; p\ln(1 + fb) + q\ln(1 - fa).
\end{equation}
By the strong law, $\frac{1}{N}\ln(W_N / W_0) \to G(f)$ almost
surely, so $W_N \approx W_0 e^{NG(f)}$ for large $N$: the
\emph{typical} wealth path grows exponentially at rate $G(f)$
\citeThorp. Unlike $\E[W_N]$, $G$ is sensitive to ruin: if
$f \geq 1/a$ there is positive probability of $W = 0$, at which
$\ln W = -\infty$ and the expectation is $-\infty$.

\subsection{First-order condition}

$G$ is strictly concave on the domain $0 < f < 1/a$ because
\begin{equation}
  G''(f) = -\frac{pb^2}{(1+fb)^2} - \frac{qa^2}{(1-fa)^2} < 0,
\end{equation}
so it has a unique interior maximum. Set $G'(f) = 0$:
\begin{align}
  G'(f) &= \frac{pb}{1+fb} - \frac{qa}{1-fa} = 0
  \\
  \Longleftrightarrow\quad pb(1-fa) &= qa(1+fb)
  \\
  \Longleftrightarrow\quad pb - pbfa &= qa + qabf
  \\
  \Longleftrightarrow\quad pb - qa &= f \cdot ab(p + q) = f\cdot ab.
\end{align}
Rearranging gives the \textbf{binary Kelly fraction}:
\begin{equation}
  \boxed{\;f^{*} \;=\; \frac{pb - qa}{ab}.\;}
  \label{eq:kelly_binary}
\end{equation}

Bet the ratio of edge to odds. The numerator $pb - qa$ is expected
profit per unit staked; the denominator $ab$ measures how violent
the outcomes are.

\subsection{Worked example: the 55/45 coin}
\label{sec:kelly_worked_55_45}

Take $p = 0.55$, $q = 0.45$, $a = b = 1$ (win or lose a dollar per
dollar staked). Equation~\eqref{eq:kelly_binary} gives
\begin{equation}
  f^{*} = \frac{0.55 \cdot 1 - 0.45 \cdot 1}{1\cdot 1}
  = 0.10.
\end{equation}
Bet $10\%$ of your bankroll per flip. The per-bet growth rate is
\begin{equation}
  G(f^{*}) = 0.55\ln(1.10) + 0.45\ln(0.90)
  \approx 0.005008,
\end{equation}
so a typical bankroll grows $\approx 0.5\%$ per bet. Over $100$
bets the typical factor is $e^{100 \cdot 0.005008} \approx 1.65$:
the median player ends up $\approx 65\%$ richer. Betting $f = 1$
instead gives $\E[W_{100}] = 1.10^{100} \approx 13{,}780$, but
with probability $1 - 0.55^{100} \approx 1 - 10^{-26}$ the player
is already broke.

\begin{keyfactbox}[Binary Kelly]
Betting repeatedly on a positive-edge two-outcome trial with win
probability $p$, win size $b$, and loss size $a$, the unique
fraction of current wealth to stake that maximises the long-run
geometric growth rate $G(f) = p\ln(1+fb) + q\ln(1-fa)$ is
\begin{equation*}
  f^{*} = \frac{pb - qa}{ab}
  \;\; = \;\; \frac{\text{edge}}{\text{odds}}.
\end{equation*}
For a fair-odds coin ($a=b=1$) this collapses to $f^{*} = p - q$,
the probability gap.
\end{keyfactbox}

\begin{intuitionbox}
The Kelly fraction maximises the \emph{logarithm} of terminal
wealth. Maximising expected wealth sizes to $100\%$ and busts on
the first loss; maximising expected log-wealth trades upside
against ruin, because $\ln$ tends to $-\infty$ near zero and so
punishes ruin infinitely. ``Bet more $=$ more profit'' is correct
only up to $f^{*}$; past it, adding size makes the \emph{typical}
player poorer while keeping the \emph{average} player richer. The
right objective for a compounding bettor is the typical player.
\end{intuitionbox}

\section{The Kelly criterion: log-normal returns}
\label{sec:kelly_lognormal}

Real markets do not deliver two-point bets. Model the risky
asset's per-period \emph{log-return} as
\begin{equation}
  r \;\sim\; \Normal(\mu, \volat^2),
\end{equation}
with $\mu$ the expected log-return and $\volat$ its standard
deviation per period. The risk-free rate per period is $\rrf$.

Hold a portfolio with weight $f$ on the risky asset and $1-f$ in
cash. $f = 0$ is fully in cash; $f = 1$ fully invested; $f > 1$
leveraged (borrow at $\rrf$); $f < 0$ short.

\subsection{Portfolio log-return and the It\^o correction}

Let $V_t$ be portfolio value. Over one period the \emph{simple}
arithmetic return of the portfolio is
\begin{equation}
  \frac{V_{t+1} - V_t}{V_t}
  = f \cdot (e^{r} - 1) + (1 - f) \cdot (e^{\rrf} - 1),
\end{equation}
because the risky asset returns $e^{r} - 1$ in simple terms over
the period and cash returns $e^{\rrf} - 1$.
Taking logs of $V_{t+1}/V_t = 1 + fR_{\text{risky}} + (1-f)
R_{\text{cash}}$ with $R = e^r - 1$ and expanding to second order
in the small-period limit (equivalently, going to continuous
time, or taking the Taylor expansion carefully), the portfolio's
log-return over the period is
\begin{equation}
  r_\text{port} \;\approx\; \rrf + f(\mu - \rrf) - \tfrac12 f^2
  \volat^2.
  \label{eq:lognormal_port_logret}
\end{equation}
The first two terms are the naive ``weighted-average'' answer: a
linear combination of risk-free and excess drift. The third,
$-\tfrac12 f^2 \volat^2$, is the \textbf{It\^o correction}.
Geometric compounding eats volatility: $+10\%$ then $-10\%$ leaves
$0.99$, not $1.00$; the loss $0.5 \cdot (0.10)^2 = 0.5\%$ is
exactly the $\tfrac12\volat^2$ drag.

It\^o's lemma (\cref{ch:black_scholes}) applied to $\ln V_t$
delivers the full derivation: the self-financing SDE
$\mathrm dV_t/V_t = (\rrf + f(\mu-\rrf))\,\mathrm dt +
f\volat\,\mathrm dW_t$ gives $\mathrm d\ln V_t = (\rrf + f(\mu-\rrf)
- \tfrac12 f^2\volat^2)\,\mathrm dt + f\volat\,\mathrm dW_t$, and
integrating recovers \eqref{eq:lognormal_port_logret}.

\subsection{Long-run growth rate and the FOC}

The long-run growth rate is, as in the binary case, the
expectation of the portfolio log-return:
\begin{equation}
  G(f) \;=\; \E[r_\text{port}]
  \;=\; \rrf + f(\mu - \rrf) - \tfrac12 f^2 \volat^2.
\end{equation}
This is a concave quadratic in $f$. Its derivative is
$G'(f) = (\mu - \rrf) - f\volat^2$, which vanishes at
\begin{equation}
  \boxed{\;f^{*} \;=\; \frac{\mu - \rrf}{\volat^2}.\;}
  \label{eq:kelly_lognormal}
\end{equation}
This is the \textbf{log-normal Kelly fraction}
\citep{thorp2006kelly, chan2013algorithmic}, the formula every
industry Kelly discussion quotes or argues about.

The value at the optimum is
\begin{equation}
  G(f^{*}) \;=\; \rrf + \frac{(\mu - \rrf)^2}{2\volat^2}
  \;=\; \rrf + \frac{1}{2}\,\text{Sharpe}^2,
\end{equation}
where $\text{Sharpe} = (\mu - \rrf)/\volat$ is the
\textbf{Sharpe ratio} (excess return per unit volatility; defined
in \cref{sec:perf_metrics}). Maximum Kelly-optimal growth depends
on Sharpe alone --- not separately on $\mu$ and $\volat$. Halving
Sharpe quarters the maximum growth rate: Sharpe $2$ is four times
better than Sharpe $1$ in compounding terms, not twice.

\subsection{Worked example: the $8\%/20\%$ strategy}
\label{sec:kelly_worked_8_20}

Take $\mu - \rrf = 0.08$ per year and $\volat = 0.20$ per year
(Thorp's numbers for the S\&P~500 \citeThorp). Then
\begin{equation}
  f^{*} = \frac{0.08}{(0.20)^2} = \frac{0.08}{0.04} = 2.0.
\end{equation}
Full Kelly prescribes $200\%$ long: borrow a dollar per dollar of
equity and invest $\$2$. The implied maximum growth rate is
\begin{equation}
  G(f^{*}) = \rrf + \frac{0.08^2}{2 \cdot 0.04} = \rrf + 0.08,
\end{equation}
an extra $8\%$ on top of $\rrf$. Half Kelly ($f = 1.0$) gives
$G(f^{*}/2) = \rrf + 0.06$: exactly three-quarters of the excess
growth at half the leverage. Both numbers are verified in
\cref{sec:kelly_reference_impl}.

$f^{*} = 2$ is the single most compelling reason no professional
runs full Kelly. It is growth-maximising under the Gaussian
assumption and insane to run: $50\%$ drawdowns would recur, each
time looking like a broken strategy to whoever holds the
position. Theory and practice fork here.

\begin{keyfactbox}[Log-normal Kelly]
For a risky asset with log-return $r \sim \Normal(\mu, \volat^2)$
per period and a risk-free rate $\rrf$ per period, the fraction
of wealth in the risky asset that maximises the long-run
geometric growth rate is
\begin{equation*}
  f^{*} = \frac{\mu - \rrf}{\volat^2}.
\end{equation*}
At this optimum the growth rate is $\rrf + \tfrac12 \text{Sharpe}^2$,
i.e.\ the maximum achievable excess growth is \emph{half the
squared Sharpe ratio} of the underlying strategy. Sharpe is the
thing that compounds; it is not ``just'' a performance metric.
\end{keyfactbox}

\subsection{Half of Kelly is half-Kelly, not half the stakes}

``Half Kelly'' means $f = f^{*}/2$, not ``bet half as often'' or
``half the dollar amount''. Kelly outputs a \emph{fraction of
current wealth}; halving it still rebalances period by period.
Compounding discipline is preserved; only the leverage knob is
turned down.

\section{Fractional Kelly and why no one trades full Kelly}
\label{sec:fractional_kelly}

The FOC says full Kelly maximises long-run growth --- in the
limit, over infinite trials, assuming $\mu$ and $\volat$ are known
exactly. None of these hold in practice, and the short-term
variance of full Kelly is brutal: $50\%$ drawdowns are routine.
Three objections explain why nobody trades it \citeThorp, \citeChan.

\subsection{Drawdown probability}

Thorp gives the following result for a log-normal portfolio at
leverage $c f^{*}$ (``$c$-Kelly''), derived as a two-barrier
Brownian-motion problem under continuous rebalancing \citeThorp:
\begin{equation}
  \Prob\!\bigl(\,\exists\, t\geq 0\text{ such that }
  V_t/V_0 \leq x\,\bigr) \;=\; x^{2/c - 1},
  \qquad 0 < x < 1.
  \label{eq:thorp_drawdown}
\end{equation}
Read: probability the portfolio is ever at or below fraction $x$
of starting value over an infinite horizon. Three cases:
\begin{itemize}
  \item \textbf{Full Kelly}, $c = 1$: exponent $1$, so
  $\Prob(\text{ever} \leq x V_0) = x$. Infinite-horizon full-Kelly
  has a $50\%$ chance of ever halving, $25\%$ of ever quartering.
  \item \textbf{Half Kelly}, $c = 1/2$: exponent $3$,
  $\Prob(\text{ever halve}) = 0.5^3 = 0.125$.
  \item \textbf{Quarter Kelly}, $c = 1/4$: exponent $7$,
  $\Prob(\text{ever halve}) = 0.5^7 \approx 0.0078$, under $1\%$.
\end{itemize}
The growth cost is much smaller than the drawdown benefit. Half
Kelly retains $3/4$ of the full-Kelly growth rate
(\cref{sec:kelly_worked_8_20}); its growth standard deviation is
exactly \emph{half} full Kelly's, variance one quarter
\citeThorp. Halving $f$ gives up $25\%$ of growth for $75\%$ less
drawdown probability --- one of the best trades in quant finance.

\subsection{Estimation error}
\label{sec:kelly_estimation_error}

$f^{*} = (\mu - \rrf)/\volat^2$ depends on $\mu - \rrf$, notoriously
hard to estimate: years of data give confidence intervals
comparable in size to $\mu$ itself. Over-estimating $\mu$
over-estimates $f^{*}$ and pushes the bettor past the peak of $G$
onto the destructive right side.

The asymmetry: the $G(f)$ curve is approximately parabolic,
peaking at $f^{*}$ and falling quadratically on both sides. Under-
betting captures positive growth, just less of it. Over-betting
far enough drives $G(f)$ \emph{negative}; typical wealth shrinks
to zero. The growth rate crosses zero at $f_c = 2f^{*}$ in the
log-normal case. Past there is terminal.

Thorp's scenario: $\hat\mu = 1.5\mu$ and $\hat f^{*} = 1.5
f^{*}_{\text{true}}$ loses about $50\%$ of typical growth
\citeThorp. $\hat\mu = 0.5\mu$ (half Kelly) retains $75\%$ of
growth. Over-betting compounds through a loss-inducing variance
drag; under-betting only loses growth.

\begin{pitfallbox}[Full Kelly is theoretically optimal and
practically insane]
Three reasons nobody runs full Kelly \citeThorp:
\begin{enumerate}
  \item \textbf{Drawdown math}. Full Kelly has a $50\%$ lifetime
  probability of at least halving and a $25\%$ chance of
  quartering. Half Kelly drops these to $12.5\%$ and $1.6\%$ at
  only a $25\%$ reduction in growth rate. The trade is lopsided in
  favour of the more conservative sizing.
  \item \textbf{Estimation error}. $f^{*}$ depends on an unknown
  $\mu$. Over-estimating $\mu$ by $50\%$ pushes you $50\%$ past
  the Kelly peak, where the growth rate collapses. The penalty
  for over-estimating is much worse than the penalty for
  under-estimating, because $G(f)$ becomes \emph{negative} at
  $f_c = 2 f^{*}$.
  \item \textbf{Behavioural intolerance}. A fund manager who
  routinely experiences $30$--$50\%$ drawdowns does not keep
  her clients, her job, or her emotional equilibrium, regardless
  of what the math says about the long run. Human tolerance for
  drawdowns is a binding constraint.
\end{enumerate}
Standard response \citeThorp: run leverage $c f^{*}$ with $c \in
[0.25, 0.5]$ and treat full Kelly as an upper bound, never a
target. Half Kelly is the default: $25\%$ growth reduction for
$75\%$ drawdown reduction, and robustness to $\mu$-estimation
error.
\end{pitfallbox}

\begin{keyfactbox}[Fractional Kelly rule of thumb]
\emph{Trade half Kelly.} Take $f = \tfrac12 f^{*}$ with
$f^{*} = (\mu - \rrf)/\volat^2$ estimated from your best data.
This halves volatility and quarters variance at $25\%$ of the
growth rate, and is robust to $\mu$ over-estimates up to a factor
of two (where full Kelly would have ruined you). Quarter Kelly
suits fat-tailed returns or an impatient client base.
\end{keyfactbox}

\begin{intuitionbox}
Fractional Kelly's appeal is not timidity. The $G(f)$ curve is
symmetric near its peak but its \emph{consequences} are
asymmetric: under-betting loses linear growth, over-betting
compounds you to zero. When $f^{*}$ cannot be measured
perfectly --- and it cannot --- the rational response is to
deliberately under-bet. Half Kelly is the optimal robust response
to estimation error in $\mu$.
\end{intuitionbox}

\section{Performance metrics}
\label{sec:perf_metrics}

Sizing is half the problem; scoring is the other. Four numbers
dominate the conversation: Sharpe, Sortino, Calmar, and
information ratio \citeChan. Formulas and failure modes follow.

\subsection{The Sharpe ratio}

Let $\{r_t\}_{t=1}^T$ be strategy returns at some frequency and
$\rrf$ the risk-free return at that frequency. The
\textbf{Sharpe ratio} is
\begin{equation}
  \text{Sharpe} \;=\;
  \frac{\bar r - \rrf}{\hat\volat_r},
  \qquad
  \bar r = \frac{1}{T}\sum_{t=1}^T r_t,
  \qquad
  \hat\volat_r^2 = \frac{1}{T-1}\sum_{t=1}^T (r_t - \bar r)^2.
\end{equation}
It is dimensionless. To compare across horizons, Sharpe is
reported \emph{annualised}:
\begin{equation}
  \text{Sharpe}_{\text{annual}} \;=\;
  \text{Sharpe}_{\text{period}} \cdot \sqrt{\text{periods per year}}.
\end{equation}
Daily annualises by $\sqrt{252} \approx 15.87$; hourly by
$\sqrt{252 \cdot 6.5}$ (US equities); monthly by $\sqrt{12}
\approx 3.46$. The $\sqrt{\cdot}$ comes from i.i.d.\ $T$-period
variance growing linearly in $T$ while the mean also grows
linearly, so their ratio scales as $T/\sqrt{T} = \sqrt{T}$.

\subsection{Worked example: a $0.04\% / 0.8\%$ daily strategy}

A strategy has $\bar r = 0.04\%$ daily and $\hat\volat_r = 0.8\%$
daily, with $\rrf = 0$. Then
\begin{equation}
  \text{Sharpe}_{\text{daily}}
  = \frac{0.0004}{0.008} = 0.05,
\end{equation}
and
\begin{equation}
  \text{Sharpe}_{\text{annual}}
  = 0.05 \cdot \sqrt{252}
  \approx 0.7937.
\end{equation}
At $\approx 0.79$ this is below the $1.0$ threshold most
systematic desks require for capital allocation. (Context:
S\&P~500 long-run Sharpe $\approx 0.3$--$0.5$; a good
market-making desk targets $2$--$5$; HFT shops double digits.)
Verified in \cref{sec:kelly_reference_impl}.

\subsection{Why Sharpe is not enough: Sortino, Calmar, IR}

Sharpe treats upside and downside symmetrically: $+5\%$
contributes as much to $\hat\volat$ as $-5\%$. For most traders
only downside is risk; upside is opportunity. The \textbf{Sortino
ratio} replaces standard deviation with \emph{downside deviation}:
\begin{equation}
  \text{Sortino}
  \;=\;
  \frac{\bar r - \rrf}{\hat\volat^{-}_r},
  \qquad
  (\hat\volat^{-}_r)^2
  \;=\; \frac{1}{T}\sum_{t=1}^T \min(r_t - \tau, 0)^2,
\end{equation}
where $\tau$ is a minimum acceptable return (typically $0$ or
$\rrf$). For symmetric returns, downside deviation $\approx
\volat/\sqrt{2}$ so Sortino $\approx \sqrt{2} \cdot$ Sharpe.
Right-skewed returns (rare big wins, frequent small losses) give
Sortino $>$ Sharpe; left-skewed (the classic ``selling
volatility'' profile) gives Sortino $<$ Sharpe, a more honest
measure.

The \textbf{Calmar ratio} uses max drawdown
(\cref{sec:drawdown_math}) as the risk measure:
\begin{equation}
  \text{Calmar}
  \;=\; \frac{\bar r_{\text{annual}}}{|\text{MaxDD}|}.
\end{equation}
Calmar $1$ means annual return equals worst drawdown; $2$ means a
year of return buys two worst drawdowns. Conservative managers
target Calmar $\geq 0.5$ and retire strategies below $0.25$.

The \textbf{information ratio (IR)} generalises Sharpe to a
benchmark $r^{b}_t$ (e.g.\ S\&P~500, or a factor-mimicking
portfolio):
\begin{equation}
  \text{IR} \;=\;
  \frac{\overline{r_t - r^{b}_t}}{\hat\volat_{r - r^{b}}}.
\end{equation}
Numerator: active return (alpha). Denominator: tracking error. An
index fund tracking perfectly has IR undefined; beating a
benchmark by $2\%$ at $4\%$ tracking error gives IR $= 0.5$.
Long-only equity managers are judged on IR because Sharpe is
dominated by the benchmark's Sharpe.

\begin{keyfactbox}[Performance metrics taxonomy]
\begin{center}
\begin{tabular}{lll}
  \toprule
  Metric & Formula & Risk measure in denominator \\
  \midrule
  Sharpe   & $(\bar r - \rrf)/\hat\volat_r$        & Full std.\ dev.\\
  Sortino  & $(\bar r - \rrf)/\hat\volat^{-}_r$   & Downside std.\\
  Calmar   & $\bar r_{\text{ann}}/|\text{MaxDD}|$   & Max drawdown \\
  IR       & $\overline{(r - r^{b})}/\hat\volat_{r-r^{b}}$
           & Tracking error \\
  \bottomrule
\end{tabular}
\end{center}
All four are scale-invariant (degree zero in wealth), annualised,
and pure numbers. Benchmarks: Sharpe $\geq 1$ to go live,
$\geq 2$ high-quality, $\geq 3$ small-capacity HFT. Sortino
tracks Sharpe within $\sqrt{2}$ for symmetric returns. Calmar
$\geq 0.5$ for capital allocation.
\end{keyfactbox}

\begin{pitfallbox}[Sharpe over $10$ is (almost always) a bug]
Annualised Sharpe above $10$ on a backtest is overwhelmingly more
likely a look-ahead bug, survivorship-biased universe, or
infeasible execution than a real edge. Real edges leak through
costs, capacity, and regime changes; they do not show Sharpe
$= 15$ net of fees. Bisect the code before celebrating (see
\cref{ch:backtesting}). The production-desk corollary is in
\cref{sec:kelly_gs}: Sharpe \emph{above} the target range gets
flagged as a bug by risk, not a victory.
\end{pitfallbox}

\section{Drawdown arithmetic}
\label{sec:drawdown_math}

Let $\{W_t\}_{t=1}^T$ be the wealth curve of a strategy. The
\textbf{running peak} at time $t$ is $P_t = \max_{s \leq t} W_s$,
the high-water mark to date. The \textbf{drawdown} at time $t$ is
\begin{equation}
  D_t \;\defeq\; \frac{W_t - P_t}{P_t} \;\leq\; 0,
\end{equation}
expressed as a negative fraction of the running peak (sometimes
as a positive percentage with the sign flipped). The
\textbf{maximum drawdown} over the sample is
\begin{equation}
  \text{MaxDD} \;\defeq\; \min_{t} D_t,
\end{equation}
the worst peak-to-trough decline ever experienced.

\subsection{The recovery asymmetry}

The key identity: to recover from a drawdown of fraction $X \in
(0,1)$ --- wealth level $(1-X)P_t$ --- the post-drawdown gain
required is $X/(1-X)$:
\begin{equation}
  (1 - X)P_t \cdot (1 + g) = P_t
  \quad\Longleftrightarrow\quad
  g = \frac{X}{1 - X}.
\end{equation}
$10\%$ requires $11.11\%$; $20\% \to 25\%$; $30\% \to 42.86\%$;
$50\% \to 100\%$; $80\% \to 400\%$. For small $X$ the gain is
$X + X^2 + \ldots$, but by $X = 0.8$ the required recovery is
\emph{five times} the drawdown.

\begin{keyfactbox}[Drawdown recovery asymmetry]
A drawdown of fraction $X$ of peak wealth requires a subsequent
gain of
\[
  g_{\text{recover}}(X) = \frac{X}{1 - X}
\]
to restore the peak. The map $X \mapsto X/(1-X)$ is linear for
small $X$, passes through $(0.5, 1.0)$, and diverges as $X \to 1$.
\begin{center}
\begin{tabular}{rc}
\toprule
Drawdown $X$ & Gain needed $X/(1-X)$ \\
\midrule
$10\%$  & $11.11\%$  \\
$20\%$  & $25.00\%$  \\
$30\%$  & $42.86\%$  \\
$50\%$  & $100\%$   \\
$80\%$  & $400\%$   \\
$90\%$  & $900\%$   \\
\bottomrule
\end{tabular}
\end{center}
\end{keyfactbox}

\subsection{Worked example: recovery time at $10\%$ annual}

A strategy compounds at $10\%$ per year and suffers a $20\%$
drawdown; required recovery $25\%$. At continuous $\mu = 0.10$,
\begin{equation}
  T_{\text{recover}}
  = \frac{\ln(1 / 0.8)}{0.10}
  = \frac{0.2231}{0.10}
  \approx 2.23~\text{years}.
\end{equation}
Two years compounding at $10\%$ just to stand still. Discrete
annual compounding gives $\log(1/0.8)/\log(1.1) \approx 2.34$
years. One mediocre year costs two good ones.

\begin{intuitionbox}
Drawdown asymmetry is why avoiding large losses matters more than
capturing large wins. Symmetry fails because of the denominator:
after losing $X$, the compounding base is $(1-X)$, not $1$.
Future percentage returns compound off a smaller base; the deeper
the hole, the disproportionately harder the climb. Half of risk
management --- stops, sizing, Kelly fractionalisation --- is
about avoiding the scenario where you need $100\%$ to break even.
\end{intuitionbox}

\section{VaR and Expected Shortfall}
\label{sec:var_es}

Two industry-standard risk metrics round out the chapter; see
\citeHull{} (Ch.\ 21) for the full treatment.

\subsection{Value at Risk}

\textbf{Value at Risk (VaR)} at confidence level $\alpha$ over
horizon $h$ is the smallest loss $L$ such that
\begin{equation}
  \Prob(\text{loss over } h \geq L) \leq 1 - \alpha.
\end{equation}
``Our $1$-day $95\%$ VaR is \$5m'' means we expect to lose more
than \$5m on at most $5\%$ of days. VaR is the
$(1-\alpha)$-quantile of the loss distribution. Typical choices:
$\alpha = 95\%$ or $99\%$; $h =$ one day (trading books) or ten
days (Basel~III capital).

VaR is popular because it reports one dollar number a
non-quantitative manager understands. It also lies about the
tail: two portfolios can share a $99\%$ VaR of \$5m but differ
wildly beyond it --- one loses exactly \$5m on its worst $1\%$
of days, the other loses \$5m on most and \$100m on a few. Same
VaR, very different behaviour when the tail fires.

\subsection{Expected Shortfall}

\textbf{Expected Shortfall (ES)}, also called
\emph{Conditional VaR} (CVaR), fixes this by averaging over the
tail:
\begin{equation}
  \text{ES}_\alpha \;=\;
  \E\!\bigl[\,\text{loss} \mid \text{loss} > \text{VaR}_\alpha\bigr].
\end{equation}
The average loss given the loss exceeds VaR. By construction
$\geq$ VaR, and it cares about tail shape, not just the $99$th
percentile. ES is \emph{coherent} (Artzner et al.): ES of a sum
of portfolios is at most the sum of individual ES's
(subadditivity). VaR fails subadditivity --- two portfolios can
have combined VaR larger than the sum of their individual VaRs,
contradicting the intuition that diversification reduces risk.
ES does not. Post-2008 Basel moved from VaR to ES for
trading-book capital \citeHull.

\begin{pitfallbox}[VaR lies about the tail]
VaR reports the \emph{boundary} of the bad region, not its
expected severity. Two portfolios can share a VaR: one just
touching the threshold, one blowing through it. Always request
Expected Shortfall alongside VaR; treat VaR alone as a compliance
box-tick. Guideline: $\text{ES}/\text{VaR} \approx 1$ means a
well-behaved tail; $\geq 1.5$ means fat-tail trouble beyond the
quantile.
\end{pitfallbox}

\begin{prosperitybox}
\textbf{PICNIC\_BASKET stat arb and half Kelly.} The Frankfurt
Hedgehogs placed 2nd in Prosperity~3 (2025) using ETF/constituent
stat-arb on PICNIC\_BASKET1/2; the basket strategy alone earned
$40$--$60$k SeaShells per round, plus related-product profits
\citep{frankfurthedgehogs2025}. Treat this as an $\{r_t\}$ series
with one ``period'' per round.

Simulator calibration for a similar Prosperity~4 strategy:
\begin{itemize}
  \item Expected PnL per round: $\hat\mu = 50{,}000$ SeaShells.
  \item Standard deviation of per-round PnL: $\hat\volat =
  10{,}000$ SeaShells (conservative, reflecting the
  mean-reverting nature of the spread and the occasional fat
  tail when the signal misfires).
  \item Starting bankroll: $W_0 = 500{,}000$ SeaShells
  (representative end-of-Round-1 bankroll for a top team).
\end{itemize}
Per-round $\mu = 50{,}000/500{,}000 = 10\%$ and $\volat =
10{,}000/500{,}000 = 2\%$ (no competition borrowing cost, so
$\rrf = 0$). Per-round Sharpe $= 0.10/0.02 = 5$. Annualising is
meaningless here --- the ``year'' is a week of rounds --- but
the per-round $5$ drives sizing.

Log-normal Kelly:
\begin{equation}
  f^{*} = \frac{\mu}{\volat^2} = \frac{0.10}{(0.02)^2}
  = \frac{0.10}{0.0004} = 250.
\end{equation}
$250\times$ leverage is nonsensical. The full-Kelly output is
telling you the competition's \emph{position limit} binds long
before Kelly does --- good news, because you should trade to
\emph{the position limit every round}, not to an abstract
``Kelly fraction''. The position limit ($60$ baskets of
PICNIC\_BASKET1) is the binding constraint. Even unconstrained,
half Kelly ($125\times$) and quarter Kelly ($62.5\times$) are
still insane.

The formula is not wrong: fractional Kelly applies when
\emph{variance is binding}. In a competition with hard position
limits the binding constraint is liquidity, not bankruptcy risk.
Applied to the $8\%/20\%$ example of
\cref{sec:kelly_worked_8_20}, where variance \emph{does} bind,
half Kelly prescribes $f = 1.0$ (fully invested). Compare to the
top-team IV-scalping VOLCANIC\_ROCK\_VOUCHERS numbers in
\cref{ch:top_team_lessons}, where vol is higher and Kelly is a
meaningful knob.
\end{prosperitybox}

\begin{gsbox}
\label{sec:kelly_gs}
On a Goldman systematic trading desk, every strategy with live
capital has a written risk profile with three numbers:
\begin{itemize}
  \item \textbf{Nominal Sharpe target}. Gross Sharpe $\geq 1.5$;
  net $\geq 1.0$. ``High-quality'' books target $\geq 2$ net.
  Below the gross target is probation; above the top flagged by
  risk as mismeasured (see the pitfall above).
  \item \textbf{Max drawdown cutoff}. Typically $10$--$15\%$,
  sometimes $5\%$ for a core market-making book. Beyond the
  cutoff the book is automatically reduced or halted: the Strat
  writes the halt logic, the trader retains override, both own
  the outcome.
  \item \textbf{Kelly fraction cap}, never above $0.5$: Strats
  compute a rolling Kelly from recent $\mu$, $\volat$, and the
  book's leverage is capped at half of that. Quarter Kelly is
  not unusual on fat-tailed books.
\end{itemize}
Strats own the live metric pipeline (rolling Sharpe, drawdown,
Kelly, alerts). Traders own the override: they can reduce a book
faster than the automated cutoff but cannot silently \emph{raise}
a Kelly cap. Separation of concerns is non-negotiable.

Two surprises for newcomers. First, the gross Sharpe target is
not much higher than the net one because desk cost models
($2$--$10$ bps per round-turn) are well-calibrated. Second,
production risk code is simpler than the research it disciplines:
Sharpe, drawdown, and Kelly are three fifty-line functions on a
five-second cadence, not a $500$-line ML pipeline. Complexity is
the enemy of live risk. See \cref{ch:risk_production} for the
data ingestion feeding these metrics.
\end{gsbox}

\begin{historybox}
\textbf{Kelly, Shannon, Thorp.} John Larry Kelly Jr.\ (Bell Labs)
derived the criterion in a 1956 paper in the \emph{Bell System
Technical Journal} \citep{kelly1956new}. His setup was
information-theoretic: how much of your bankroll to bet on a
sequence of noisy tips. The growth-optimal answer maximises
$\E[\ln W]$, and the quantity inside turned out to be the mutual
information between tips and outcome --- a surprising link from
gambling to Shannon's theory.

Claude Shannon handed the paper to Edward~O. Thorp at MIT in
1960, just after Thorp had solved blackjack card counting.
Thorp's \emph{Beat the Dealer} (1962) applied Kelly to blackjack;
he co-founded Princeton--Newport Partners in 1969. Over its
19-year life, Princeton--Newport compounded at $\approx 19.1\%$
net of fees, never had a losing year, and used fractional Kelly
throughout \citeThorp. Nobody there traded full Kelly.
\end{historybox}

\section{A short reference implementation}
\label{sec:kelly_reference_impl}

\begin{lstlisting}[language=Python, caption={Kelly and
performance-metric reference implementation.}]
import numpy as np

def kelly_binary(p, b, a):
    """Optimal fraction for a two-outcome bet. Win b with prob p, lose a with prob 1-p."""
    return (p * b - (1 - p) * a) / (a * b)

def kelly_lognormal(mu, sigma, r_f=0.0):
    """Optimal fraction for a log-normal risky asset vs cash at r_f."""
    return (mu - r_f) / sigma ** 2

def sharpe(returns, r_f=0.0, periods_per_year=252):
    """Annualised Sharpe from a per-period return series."""
    excess = returns - r_f / periods_per_year
    return excess.mean() / excess.std(ddof=1) * np.sqrt(periods_per_year)

def max_drawdown(wealth):
    """Worst peak-to-trough decline on a wealth curve (fraction, <= 0)."""
    peak = np.maximum.accumulate(wealth)
    return ((wealth - peak) / peak).min()

# Sanity checks (verified numerically in chapter text):
assert abs(kelly_binary(0.55, 1, 1) - 0.10) < 1e-12
assert abs(kelly_lognormal(0.08, 0.20) - 2.0) < 1e-12
\end{lstlisting}

These four functions implement the chapter's sizing and scoring.
A production library adds rolling windows, downside deviation,
drawdown recovery time, IR, and tail-conditioned loss --- each a
wrapper around these primitives.

\section{Key facts to memorise}
\label{sec:kelly_key_facts}

\begin{itemize}
  \item \textbf{Binary Kelly}: $f^{*} = (pb - qa)/(ab)$. For a
  fair-odds coin, $f^{*} = p - q$. For the $55/45$ coin, bet
  $10\%$ per trial.
  \item \textbf{Log-normal Kelly}: $f^{*} = (\mu - \rrf)/\volat^2$.
  Maximum achievable long-run growth is $\rrf + \tfrac12
  \text{Sharpe}^2$. Sharpe squared is what compounds.
  \item \textbf{Fractional Kelly}: trade $cf^{*}$ with $c \in
  [0.25, 0.5]$. Half Kelly retains $75\%$ of the growth rate
  with half the volatility and a quarter of the variance;
  Thorp's rule of thumb.
  \item \textbf{Drawdown probability at $c$-Kelly}: the lifetime
  probability of ever drawing down to fraction $x$ of starting
  wealth is $x^{2/c - 1}$. Full Kelly: $50\%$ chance of ever
  halving. Half Kelly: $12.5\%$. Quarter Kelly: $0.78\%$.
  \item \textbf{Sharpe}: $(\bar r - \rrf)/\hat\volat_r$,
  annualised with $\sqrt{252}$ from daily, $\sqrt{12}$ from
  monthly. Gross $\geq 1.5$, net $\geq 1.0$ is a live-capital
  threshold on a systematic desk. $0.79$ (the worked example) is
  sub-threshold.
  \item \textbf{Sortino}: use downside deviation in the
  denominator. Pay attention to it when returns are skewed.
  \item \textbf{Calmar}: annual return divided by $|$max
  drawdown$|$. A direct tradeoff of growth against worst
  historical loss. Conservative managers target $\geq 0.5$.
  \item \textbf{Information ratio}: alpha over tracking error.
  Right metric for a long-only or benchmarked book.
  \item \textbf{Drawdown recovery asymmetry}: loss of fraction
  $X$ requires gain of $X/(1-X)$ to recover. $20\% \to 25\%$;
  $50\% \to 100\%$; $80\% \to 400\%$. This is the single most
  important identity in practical risk management and it is why
  avoiding large losses beats capturing large wins.
  \item \textbf{VaR vs ES}: VaR is the $(1-\alpha)$-quantile of
  loss; ES is the \emph{expected} loss beyond VaR. ES is
  coherent (subadditive), VaR is not. Post-Basel~III the
  regulatory standard is ES.
  \item \textbf{Kelly cap in production}: strats run rolling
  Kelly estimates; the production book is capped at half of the
  Kelly estimate, never at full Kelly. Sizing is half the risk
  problem; scoring (rolling Sharpe, drawdown, Kelly, VaR, ES) is
  the other half.
\end{itemize}
